\chapter{Why Study Microeconomics?}
\label{chap:introduction}

\section{Learning objectives}

After completing this chapter, you should be able to:

\begin{itemize}
  \item explain what microeconomics studies and distinguish it from macroeconomics;
  \item identify scarcity, opportunity cost, incentives, and institutions in everyday decisions;
  \item distinguish positive from normative economic claims;
  \item explain why economists use models and assess the role of their assumptions; and
  \item relate microeconomic reasoning to central questions in politics, philosophy, and economics.
\end{itemize}

\section{Economics begins with choice}

Every individual, firm, and government faces choices. A student has limited time to allocate between studying, working, resting, and social life. A firm must decide which products to offer, which workers and machines to employ, and what price to charge. A government must decide how to allocate public resources among health care, education, defence, infrastructure, and many other priorities.

These choices arise because resources are scarce. 
\keyterm{Scarcity} means that the resources available to satisfy human wants are limited. Time, income, labour, land, natural resources, and information are all limited. Scarcity does not mean that a person is poor; even a very wealthy person has only twenty-four hours in a day.

Economics studies how people and institutions make choices under scarcity, and how those choices interact. It does not assume that people care only about money. People may value friendship, fairness, leisure, political influence, environmental quality, or the well-being of others. The economic question is nevertheless the same: given limited resources, which option is chosen, and what is given up?

\section{What is microeconomics?}

\keyterm{Microeconomics} studies the choices of individual economic agents and the way those choices are coordinated in markets and other institutions. Its basic units of analysis are consumers, households, workers, firms, and particular markets.

Some typical microeconomic questions are:

\begin{itemize}
  \item How does a consumer divide a limited budget between food, housing, transport, and leisure?
  \item How does a firm choose its output when producing more requires additional workers or machinery?
  \item Why do some markets have many small firms while others are dominated by one or a few firms?
  \item What happens to consumers and producers when a government introduces a tax, a subsidy, or a price ceiling?
  \item When do private decisions create costs or benefits for people who are not directly involved in the decision?
\end{itemize}

Microeconomics differs from \keyterm{macroeconomics}, which studies economy-wide aggregates such as inflation, unemployment, economic growth, interest rates, and national income. The distinction is useful, but it is not absolute. Aggregate outcomes ultimately depend on many individual decisions, while individual decisions are affected by the wider macroeconomic environment.

\section{Opportunity cost}

Making a choice means rejecting an alternative. The value of the best alternative forgone is the \keyterm{opportunity cost} of a decision.

Suppose that you have two hours available in the afternoon. You can revise for an exam, work in a paid job, or meet friends. If you choose to revise and the best alternative was the paid job, the income you could have earned is part of the opportunity cost of studying. If the best alternative was meeting friends, the lost leisure is part of the opportunity cost instead.

\begin{figure}[htbp]
  \centering
  \begin{tikzpicture}[x=0.88cm,y=0.88cm]
    \draw[->,PPEInk,line width=0.7pt] (0,0) -- (6.2,0) node[below right,font=\sffamily\small]{Education services};
    \draw[->,PPEInk,line width=0.7pt] (0,0) -- (0,4.8) node[above left,font=\sffamily\small]{Health-care services};
    \draw[PPENavarraRed,line width=1.15pt] (0,4.25) .. controls (2.3,4.05) and (4.75,2.65) .. (5.75,0);
    \fill[PPEBurgundy] (3.05,3.32) circle (2.2pt) node[above right,font=\sffamily\small]{E: efficient};
    \fill[PPEGold] (2.0,1.65) circle (2.2pt) node[below left,font=\sffamily\small]{I: inefficient};
    \fill[PPENavarraRed] (4.55,3.62) circle (2.2pt) node[above right,font=\sffamily\small]{U: unattainable};
    \draw[PPEGold,dashed] (3.05,0) -- (3.05,3.32) -- (0,3.32);
    \draw[ppearrow] (2.38,3.73) to[bend right=16] node[above,font=\sffamily\scriptsize,text=PPEBurgundy]{opportunity cost} (3.68,2.90);
  \end{tikzpicture}
  \caption{Scarcity makes some combinations feasible, some unattainable, and some inefficient. Moving along the frontier requires a trade-off.}
  \label{fig:scarcity-ppf}
\end{figure}

Opportunity cost is broader than a monetary payment. Attending university involves tuition and living expenses, but it may also involve forgone earnings, forgone leisure, and delayed entry into the labour market. Similarly, a government programme uses not only money but also public attention, administrative capacity, and political capital.

\section{Incentives and marginal thinking}

People respond to \keyterm{incentives}: changes in the costs or benefits associated with an action. A higher wage may encourage someone to work more hours. A tax on cigarettes may reduce cigarette consumption. A subsidy for vaccination may increase the number of people vaccinated.

Economic reasoning often focuses on small changes rather than all-or-nothing choices. A student deciding whether to study one additional hour compares the expected benefit of that additional hour with its cost in lost leisure, work, or sleep. This is called \keyterm{marginal thinking}. A rational decision-maker continues an activity while its marginal benefit is at least as large as its marginal cost.

Marginal thinking does not imply that people always calculate precisely. Rather, it provides a disciplined way to identify the relevant trade-off. Later chapters will express this idea formally through marginal utility, marginal product, marginal cost, and marginal revenue.

\begin{keyidea}
\textbf{A discrete marginal decision.} Suppose that a student's additional study hours have the following estimated benefits and costs, expressed in euro equivalents.
\begin{center}
\begin{tabular}{lrr}
\toprule
Additional study hour & Marginal benefit & Marginal cost \\
\midrule
1 & \euro8 & \euro2 \\
2 & \euro6 & \euro3 \\
3 & \euro4 & \euro4 \\
4 & \euro2 & \euro5 \\
\bottomrule
\end{tabular}
\end{center}
The first three hours are worthwhile in this simple example; the fourth is not. For a continuous choice with a well-behaved objective, an interior optimum is normally characterised by marginal benefit equal to marginal cost.
\end{keyidea}

\section{Rationality and \emph{ceteris paribus}}

Economists often model an agent as choosing the feasible option that best advances that agent's objectives. This is the \keyterm{rational-choice assumption}. It does not mean that people are selfish, perfectly informed, or able to calculate without error. A student may care about friendship or fairness; a firm may care about reputation; and a voter may care about principle. Rationality means only that the model treats choices as purposeful and constrained.

\keyterm{Ceteris paribus} means ``with other relevant factors held constant.'' It allows economists to isolate one relationship at a time. For example, the law of demand asks how quantity demanded changes when price changes, holding income, tastes, and other determinants fixed. When this simplifying assumption is not adequate, later chapters consider strategic interaction, incomplete information, and behavioural economics.

\section{Institutions matter}

Choices are made within rules. \keyterm{Institutions} are the formal and informal rules that structure social interaction: laws, property rights, contracts, courts, social norms, political procedures, and market arrangements.

The same resources can generate very different outcomes under different institutions. Secure property rights may encourage investment because people expect to benefit from the returns to their effort. Clear rules of contract may make exchange easier between people who do not know each other. Conversely, weak enforcement or corruption can discourage trade and productive investment.

For PPE students, this point is particularly important. Microeconomics is not only a study of prices and firms; it is also a framework for examining how rules shape incentives, power, cooperation, and distribution.

\begin{exercisebox}
\textbf{Running example: changing the assessment rule.} Suppose a degree programme changes a module from one final examination to $60\%$ continuous assessment and $40\%$ final examination. Students still have the same limited time, but the new rule changes the return to attendance, weekly preparation, group work, and last-minute revision. It may also change which students perform best.

The rule does not determine every student's behaviour. It changes the incentives within which students choose, and those choices shape the final distribution of marks. This sequence---rules, incentives, choices, and allocations---is one of the book's central themes.
\end{exercisebox}

\begin{figure}[htbp]
  \centering
  \begin{tikzpicture}[node distance=0.32cm]
    \node[ppebox,text width=2.65cm] (agent) {Decision-maker};
    \node[ppebox,text width=2.65cm,right=of agent] (alternatives) {Feasible alternatives};
    \node[ppebox,text width=2.65cm,right=of alternatives] (incentives) {Incentives and trade-offs};
    \node[ppebox,text width=2.65cm,right=of incentives] (consequences) {Consequences for others};
    \draw[ppearrow] (agent) -- (alternatives);
    \draw[ppearrow] (alternatives) -- (incentives);
    \draw[ppearrow] (incentives) -- (consequences);
    \node[draw=PPEGold,fill=PPEGold!14,rounded corners=3pt,below=0.5cm of alternatives,minimum width=9.4cm,
      minimum height=0.8cm,align=center,font=\sffamily\small] {Institutions: rules, rights, norms, and enforcement};
  \end{tikzpicture}
  \caption{A reusable framework for microeconomic reasoning in PPE.}
  \label{fig:economic-lens}
\end{figure}

\section{Economic models and assumptions}

The social world is complex. Economists therefore use \keyterm{models}: simplified representations of reality designed to answer a specific question. A model deliberately leaves out many details so that the relevant mechanisms can be analysed clearly.

For example, a model of consumer choice may initially assume that a consumer knows prices, has a fixed income, and has stable preferences. These assumptions do not claim that real people are perfectly informed or never change their minds. They allow us to study how income and prices affect choices before adding further complications.

A useful model should be judged by the question it answers, the clarity of its logic, and the usefulness of its predictions or explanations. An unrealistic assumption is not automatically a defect. The relevant question is whether simplifying that feature prevents the model from explaining the phenomenon we wish to understand.

\section{Positive and normative economics}

\keyterm{Positive economics} describes and explains the world as it is. For example, ``a higher consumer price of petrol reduces the quantity of petrol demanded, other things equal'' is a positive claim. It can be assessed using theory and evidence. A petrol tax may produce this result to the extent that it raises the price paid by consumers.

\keyterm{Normative economics} makes a judgement about what ought to be done. For example, ``the government should increase the tax on petrol'' is a normative claim. Such a claim may use positive economic analysis, but it also depends on values: how society weighs environmental protection, income distribution, mobility, and individual freedom.

The distinction is essential in PPE. Economics can clarify the likely consequences of a policy, identify who gains and loses, and reveal trade-offs. It cannot by itself determine which distribution of costs and benefits is morally or politically desirable.

\begin{keyidea}
\begin{center}
\begin{tabular}{p{0.41\textwidth}p{0.41\textwidth}}
\textbf{\color{PPEBurgundy}Positive economics} & \textbf{\color{PPENavarraRed}Normative economics} \\
\midrule
\textit{What happens if \ldots?} & \textit{What ought to be done?} \\
Uses theory and evidence to assess claims & Requires values and ethical or political judgement \\
Example: a higher petrol price lowers quantity demanded & Example: petrol taxes should be increased \\
\end{tabular}
\end{center}
\smallskip
\textbf{Policy analysis combines predicted consequences with explicit value judgements.}
\end{keyidea}

\section{A roadmap for this book}

The book begins with consumer theory: preferences, utility, budgets, demand, and elasticity. It then turns to firms, production technology, costs, and supply. These tools are combined to analyse competitive markets, monopoly, and oligopoly. The final parts examine general equilibrium, the case for and limits of markets, market failures, strategic interaction, auctions, and behavioural economics.

The central habit throughout is simple: identify the decision-maker, the feasible alternatives, the incentives, and the consequences of the choice for others.

\section*{Chapter Review}

\subsection*{Key concepts}
\begin{itemize}
  \item \textbf{Scarcity}: limited resources relative to human wants.
  \item \textbf{Opportunity cost}: the value of the best alternative forgone.
  \item \textbf{Microeconomics}: the study of individual choices and their interaction in markets and institutions.
  \item \textbf{Incentives}: factors that alter the costs or benefits of an action.
  \item \textbf{Marginal thinking}: comparing the benefit and cost of a small additional action.
  \item \textbf{Institutions}: formal and informal rules that structure interaction.
  \item \textbf{Positive economics}: explanatory claims that can be assessed with evidence.
  \item \textbf{Normative economics}: claims about what ought to be done.
\end{itemize}

\subsection*{Key questions}
\begin{itemize}
  \item What is scarce?
  \item What are the feasible alternatives?
  \item What is the opportunity cost of each choice?
  \item How do incentives and institutions affect behaviour?
  \item Who gains and who loses from a decision or policy?
\end{itemize}

\section*{Exercises}

\subsection*{Conceptual questions}
\begin{enumerate}[label=\arabic*.]
  \item Define scarcity. Why does scarcity exist even in high-income societies?
  \item Explain the difference between microeconomics and macroeconomics. Give two examples of questions belonging to each field.
  \item Define opportunity cost. Why is it not equivalent to monetary expenditure?
  \item Explain marginal thinking using a decision to study for one additional hour.
  \item What is an economic model? Why can a model be useful even when some of its assumptions are unrealistic?
\end{enumerate}

\subsection*{Graphical exercises}
\begin{enumerate}[resume,label=\arabic*.]
  \item A university has at most $120$ teaching-assistant hours available each week. An economics support session, $x$, uses $3$ hours; a political-philosophy support session, $y$, uses $2$ hours. Treat sessions as divisible for this model.
  \begin{enumerate}[label=(\alph*)]
    \item Write the feasible set and the production possibilities frontier.
    \item Find the intercepts and draw the frontier, with $x$ on the horizontal axis and $y$ on the vertical axis.
    \item Classify $A=(20,30)$, $B=(20,15)$, and $C=(30,20)$ as efficient, feasible but inefficient, or infeasible.
    \item What is the opportunity cost of one additional economics support session in units of political-philosophy support sessions? How does it relate to the slope of the frontier?
  \end{enumerate}
  \item A city introduces a congestion charge for cars entering its centre. Identify two groups likely to be affected, explain how their incentives change, and draw a simple diagram that represents one resulting trade-off.
\end{enumerate}

\subsection*{Analytical exercises}
\begin{enumerate}[resume,label=\arabic*.]
  \item Marta has one free evening. She values attending a public-policy lecture at \euro18, working a three-hour shift at \euro36 net of effort, and meeting friends at \euro25. She attends the lecture. Identify the opportunity cost and explain why the other two values should not be added.
  \item Classify each statement as positive or normative. Briefly justify your answer.
  \textit{Classify the type of claim, not whether you believe it is empirically true.}
  \begin{enumerate}[label=(\alph*)]
    \item Increasing the minimum wage reduces employment among low-skilled workers.
    \item The government should increase the minimum wage.
    \item A tax on carbon emissions increases firms' private cost of producing carbon-intensive goods.
    \item Climate policy should give priority to protecting future generations.
  \end{enumerate}
\end{enumerate}

\subsection*{Exam-style multiple-choice questions}
\begin{enumerate}[resume,label=\arabic*.]
  \item A student declines a paid shift that would yield a net \euro45 and a concert valued at \euro30 in order to attend a voluntary revision workshop. Assuming these are the only relevant alternatives, the opportunity cost of attending the workshop is:
  \begin{enumerate}[label=(\alph*)]
    \item \euro75.
    \item \euro45.
    \item \euro30.
    \item zero, because revision is educational.
  \end{enumerate}

  \item Which statement is normative?
  \begin{enumerate}[label=(\alph*)]
    \item A subsidy lowers the price paid by consumers, other things equal.
    \item Demand for a normal good rises when income rises, other things equal.
    \item Governments should subsidise public transport.
    \item A price ceiling below equilibrium creates excess demand.
  \end{enumerate}

  \item A model is most useful when it:
  \begin{enumerate}[label=(\alph*)]
    \item includes every feature of the real world.
    \item avoids making assumptions.
    \item isolates the mechanisms relevant to the question being studied.
    \item produces conclusions that cannot be tested or discussed.
  \end{enumerate}
\end{enumerate}

\subsection*{PPE discussion question}

\noindent ``Markets are institutions, not natural facts.'' Discuss this statement. In your answer, explain how property rights, contract enforcement, and social norms can affect economic exchange.

% Future chapters will be added here.
